% !TEX program = pdflatex
% !TEX encoding = UTF-8 Unicode

% ====================================================================
% ARCHIVO PRINCIPAL — Práctica Grupal Temas 4 y 5
% Asignatura: Dirección de Recursos Humanos I
% Doble Grado ADE-Ingenierías — Universidad de Granada
% Empresa analizada: CaixaBank
%
% INSTRUCCIONES DE COMPILACIÓN:
%   pdflatex practicat4t5.tex && pdflatex practicat4t5.tex
%   (dos pasadas para generar índice y referencias correctamente)
%
% Para la bibliografía, si añadís citas al archivo library.bib:
%   pdflatex practicat4t5.tex && bibtex practicat4t5 && pdflatex practicat4t5.tex && pdflatex practicat4t5.tex
% ====================================================================

\documentclass[print, color]{ugrTFG}

% -------------------------------------------------------------------
%  VERSIÓN ELECTRÓNICA PARA TABLETA
%  Cambiad "print" por "tablet" para generar un PDF adaptado a tabletas.
% -------------------------------------------------------------------

% -------------------------------------------------------------------
%  INFORMACIÓN DEL TRABAJO Y LOS AUTORES
% -------------------------------------------------------------------

\renewcommand{\miTipoTrabajo}{Práctica Grupal}
\renewcommand{\miAsignatura}{Dirección de Recursos Humanos I}

\newcommand{\miTitulo}{Reclutamiento, Selección e Integración de Personas en CaixaBank\xspace}

% Para un trabajo grupal, \miNombre contiene todos los autores.
% En la portada aparecerá como "Presentado por:" seguido de los nombres.
% Ajustadlos con vuestros nombres reales.
\newcommand{\miNombre}{%
  José Angel Carretero Montes \\ 
  David Bacas Posadas \\
  Julian Carrion Tovar \\
  Jesús Rodriguez Gonzalez \\
  Ismael Sallami Moreno \\
}

\newcommand{\miGrado}{Doble Grado ADE-Ingenierías}
\newcommand{\miFacultad}{Facultad de Ciencias Empresariales y Escuela Técnica Superior de Ingenierías}
\newcommand{\miUniversidad}{Universidad de Granada}

% Nombre del profesor/a responsable de la asignatura
\newcommand{\miTutor}{
  Javier Aguilera Caracuel \\ \emph{Departamento de Organización de Empresas II}
}

\newcommand{\miCurso}{2025-2026}

% -------------------------------------------------------------------
%  PAQUETES ADICIONALES PARA ESTE TRABAJO
% -------------------------------------------------------------------

\usepackage{tabularx}       % Tablas con columnas de anchura automática
\usepackage{multirow}       % Celdas que ocupan varias filas
\usepackage{float}          % Control de posición de figuras [H]
\usepackage{eurosym}        % Símbolo del euro: \euro{}
\usepackage{enumitem}       % Control avanzado de listas
\usepackage{tikz}           % Gráficos vectoriales (organigrama)
\usepackage{rotating}       % Figuras apaisadas (sidewaysfigure, sin problemas twoside)

% Cajas de color para recordatorios al equipo
\usepackage[most]{tcolorbox}
\tcbuselibrary{skins, breakable}

% Color corporativo UGR
\definecolor{ugrazul}{RGB}{0,75,135}

% Entorno "notaequipo": caja amarilla para recordatorios y preguntas guía.
% Usadlo cuando queráis dejar instrucciones a vuestros compañeros sobre
% qué información hay que buscar o qué preguntar en la entrevista.
\newtcolorbox{notaequipo}{
  enhanced,
  breakable,
  colback=yellow!15!white,
  colframe=orange!70!black,
  fonttitle=\sffamily\bfseries,
  title={Recordatorio del equipo},
  left=6pt, right=6pt, top=4pt, bottom=4pt,
  before skip=8pt, after skip=8pt
}

% -------------------------------------------------------------------
%  METADATOS PDF
% -------------------------------------------------------------------
\hypersetup{
  pdftitle={\miTitulo},
  pdfauthor={Grupo --- Dirección de Recursos Humanos I --- UGR},
  pdfsubject={Práctica Grupal — CaixaBank}
}

% ====================================================================
\begin{document}
% ====================================================================

\maketitle

% -------------------------------------------------------------------
%  FRONTMATTER
%  (páginas previas numeradas en romano: resumen, índice)
% -------------------------------------------------------------------
\frontmatter

% Resumen ejecutivo del trabajo
% !TeX root = ../practicat6t7t8.tex
% !TeX encoding = utf8

\chapter{Resumen}

El presente trabajo analiza las políticas de Formación y Desarrollo
(Tema~6), Evaluación y Gestión del Rendimiento (Tema~7) y Retribución
(Tema~8) aplicadas en CaixaBank, S.A., primer grupo bancario español
por volumen de activos, con más de 45.000 empleados.

En el ámbito de la formación, se examina el ciclo completo del programa
formativo de la entidad —identificación de necesidades, diseño,
implantación y evaluación—, así como el tiempo y coste asociados.
Se analizan igualmente las nuevas tendencias adoptadas: la plataforma
de \textit{e-learning} Virtaula, el modelo 70-20-10, la estructura de
universidades corporativas (escuelas especializadas internas),
el \textit{outdoor training} y la gamificación. Respecto a la gestión
de la carrera profesional, se describe el proceso de planificación
en sus cuatro fases —premisas, valoración, dirección y desarrollo—,
detallando las herramientas de rotación de puestos, \textit{coaching},
\textit{mentoring} e \textit{mentoring} inverso empleadas por la entidad,
junto con las tendencias de \textit{networking}, empleabilidad y marca
personal.

En el ámbito de la evaluación del rendimiento, se estudia la delimitación
de la evaluación (qué, cuándo y quién evalúa), los métodos objetivos,
subjetivos y mixtos utilizados, así como las estrategias de reforzamiento
positivo, negativo y la entrevista de evaluación. Se presta especial
atención a las tendencias de evaluación continua y por compromisos que
CaixaBank está implantando en sustitución de los sistemas anuales
tradicionales.

Finalmente, en materia retributiva, se analiza la estructura salarial de
cinco puestos representativos con referencia al Convenio Colectivo
sectorial de Banca, los tres tipos de equidad (interna, externa e
individual), el modelo híbrido de retribución por puesto y por
competencias, los principales incentivos financieros y elementos no
financieros, y las tendencias de retribución flexible y salario emocional.

La metodología combina el análisis de fuentes documentales públicas
—memoria anual, informe de sostenibilidad y Plan Estratégico
2025--2027 de CaixaBank— con una entrevista semiestructurada a una
directora de sucursal de la entidad.

\medskip

\textbf{Palabras clave:} formación corporativa, \textit{e-learning},
gestión de carrera profesional, \textit{coaching}, \textit{mentoring},
evaluación del rendimiento, retribución flexible, salario emocional,
CaixaBank, convenio colectivo de banca.

\cleardoublepage
\endinput


% Tabla de contenidos
% !TeX root = ../tfg.tex
% !TeX encoding = utf8
%
% ====================================================================
%  TABLA DE CONTENIDOS
% ====================================================================

\phantomsection
\pdfbookmark[0]{\contentsname}{toc}

\setcounter{tocdepth}{2}      % El índice muestra hasta subsecciones
\setcounter{secnumdepth}{3}   % Se numeran hasta subsubsecciones

\tableofcontents

% Listas de figuras y tablas (descomentad si las necesitáis)
% \phantomsection
% \listoffigures

% \phantomsection
% \listoftables

\cleardoublepage


% -------------------------------------------------------------------
%  MAINMATTER
%  (cuerpo del trabajo, numerado en arábigo)
% -------------------------------------------------------------------
\mainmatter

% ---- Capítulo 4: Tema 4 — Reclutamiento ----
% ====================================================================
% CAPÍTULO 4 — TEMA 4: RECLUTAMIENTO
% Guía profesor:
%   1) Costes
%   2) Tiempo
%   3) Tipo (interno, externo, mixto)
%   4) Métodos (internos y externos)
%   5) Nuevas tendencias (2.0, 3.0, vídeo-CV)
%   6) Propuestas de mejora
% ====================================================================

\chapter{Tema 4: Reclutamiento}

% ---- Bloque Jesús ----
% ====================================================================
% TEMA 4 — JESÚS
% Bases y Métodos de Reclutamiento
% ====================================================================

\section{Tema 4: Bases y Métodos de Reclutamiento}

\subsection{Costes del proceso de reclutamiento}

\subsection{Tiempo empleado en el proceso de reclutamiento}

\subsection{Tipo de reclutamiento aplicado}

\subsubsection{Reclutamiento interno}

\subsubsection{Reclutamiento externo}

\subsubsection{Reclutamiento mixto}

\subsection{Métodos de reclutamiento utilizados}

\subsubsection{Métodos de reclutamiento interno}

\subsubsection{Métodos de reclutamiento externo}


% ---- Bloque José Ángel ----
% ====================================================================
% TEMA 4 — JOSÉ ÁNGEL
% Innovación y Propuestas de Reclutamiento
% ====================================================================

\section{Innovación y Propuestas de Reclutamiento}

\subsection{Nuevas tendencias de reclutamiento: modelos 2.0 y 3.0}

La digitalización ha transformado profundamente el reclutamiento. La teoría distingue el
\textbf{e-recruitment} o \textbf{reclutamiento 2.0}, basado en el uso de internet y portales
de empleo, del \textbf{reclutamiento 3.0} o \textit{social recruiting}, definido como «el
proceso por el cual las empresas buscan e identifican los posibles candidatos para cubrir
sus necesidades internas de empleo a través de las redes sociales multiplataformas de
internet» [8].

CaixaBank combina ambos modelos de forma madura:

\begin{itemize}
  \item \textbf{Reclutamiento 2.0.} La entidad dispone de un portal propio,
        \textit{CaixaBankCareers}, donde los candidatos registran su candidatura de forma
        digital, acceden a información sobre el proceso y realizan las pruebas en línea
        [1]. Además, colabora con agencias de trabajo temporal como
        Randstad para ampliar el alcance externo del proceso.

  \item \textbf{Reclutamiento 3.0.} CaixaBank mantiene una presencia activa en LinkedIn
        —con más de 195\,000 seguidores— donde publica ofertas, difunde contenido de
        \textit{employer branding} y contacta directamente con perfiles pasivos
        [5]. En 2019 lanzó el \textbf{People Xperience Hub},
        una iniciativa de \textit{Recruitment Process Outsourcing} (RPO) que estandariza
        y digitaliza todo el ciclo de atracción de talento [6].
        También utiliza Instagram y otras redes sociales para reforzar su marca de
        empleador.
\end{itemize}

Esta estrategia presenta las ventajas señaladas en la teoría: menor coste respecto al
reclutamiento tradicional, mayor alcance geográfico y mejor capacidad de precalificación
y filtrado de candidatos. No obstante, CaixaBank debe gestionar con cuidado el riesgo de
reputación \textit{online} y el cumplimiento normativo en el tratamiento de datos
personales de los candidatos.

\subsection{Uso e impacto del Vídeo-CV}

El \textbf{videocurrículum} es una tendencia consolidada en los procesos modernos de
reclutamiento: el candidato se graba presentando su perfil, permitiendo al equipo de
selección valorar competencias comunicativas y actitudinales que escapan al CV
tradicional [8].

El proceso de selección de CaixaBank consta de seis fases digitalizadas [1], lo que
lo hace compatible con la incorporación del videocurrículum como herramienta de
preselección. Su estructura online —con pruebas y dinámicas realizadas en remoto—
encaja directamente con el modelo que describe la teoría.

El impacto de esta herramienta sería doble. Para la empresa, reduce desplazamientos y
costes logísticos, acelera la criba y permite revisar las grabaciones de forma asíncrona.
Para el candidato, ofrece mayor flexibilidad horaria y elimina barreras geográficas,
aunque puede introducir sesgos si no se establecen criterios de evaluación estructurados.
Un riesgo adicional es la brecha digital: candidatos sin acceso a tecnología adecuada
pueden quedar en desventaja, lo que puede comprometer la equidad del proceso.

\subsection{Propuestas de mejora para el reclutamiento futuro}

\subsubsection{Mejoras a corto plazo}

\begin{itemize}
  \item \textbf{Incorporación de Inteligencia Artificial en el cribado.} La teoría
        recoge que herramientas de \textit{machine learning} permiten filtrar currículos
        automáticamente y realizar búsquedas semánticas para detectar candidatos idóneos
        [8]. CaixaBank podría integrar un sistema ATS (\textit{Applicant
        Tracking System}) con IA para reducir el tiempo de criba y minimizar sesgos
        inconscientes, garantizando siempre la supervisión humana de las decisiones
        finales.

  \item \textbf{Gamificación del proceso de selección.} La introducción de pruebas
        situacionales en formato de juego (\textit{serious games}) durante la fase de
        preselección permitiría evaluar competencias como la toma de decisiones o el
        trabajo en equipo de forma más objetiva y atractiva para los candidatos,
        mejorando a la vez la experiencia de candidatura (\textit{candidate experience}).
\end{itemize}

\subsubsection{Mejoras a medio plazo}

\begin{itemize}
  \item \textbf{Potenciación del \textit{employer branding} en redes sociales.}
        Aunque CaixaBank ya cuenta con presencia en LinkedIn, podría intensificar
        el uso de contenido generado por empleados (\textit{employee advocacy}) y
        la difusión de testimonios reales para atraer talento pasivo, reduciendo
        la dependencia de canales de pago y mejorando su posicionamiento como
        empleador de referencia en el sector financiero.

  \item \textbf{Programas de reclutamiento universitario (\textit{campus recruiting}).}
        La teoría señala a las instituciones docentes como fuente clave para obtener
        candidatos cualificados [8]. CaixaBank podría estructurar
        acuerdos estables con universidades —más allá de las prácticas actuales— para
        identificar talento en etapas tempranas y construir una bolsa de candidatos
        preseleccionados para posiciones de perfil junior.
\end{itemize}

\subsubsection{Indicadores de seguimiento propuestos}

Para evaluar la eficacia de las mejoras, se propone monitorizar los siguientes
indicadores, coherentes con los ratios de eficiencia del reclutamiento definidos
en la teoría [8]:

\begin{itemize}
  \item \textbf{\textit{Yield ratio}}: proporción de candidatos que superan cada
        fase del proceso, para identificar cuellos de botella.
  \item \textbf{\textit{Time-to-hire}}: tiempo medio desde la publicación de la
        oferta hasta la contratación, con el objetivo de reducirlo por debajo de
        los 41 días actuales.
  \item \textbf{Tasa de aceptación de oferta}: porcentaje de candidatos que aceptan
        la propuesta, indicador de la percepción de la marca empleadora.
  \item \textbf{Satisfacción del candidato}: encuesta post-proceso para medir la
        \textit{candidate experience} e identificar áreas de mejora.
\end{itemize}



% ---- Capítulo 5: Tema 5 — Selección e Integración ----
% ====================================================================
% CAPÍTULO 5 — TEMA 5: SELECCIÓN E INTEGRACIÓN
% Guía profesor:
%   - Etapas del proceso de selección
%   - Herramientas de selección y su gestión
%   - Nuevas tendencias (vídeo entrevistas)
%   - Socialización formal e informal
%   - Manual de acogida
%   - Propuestas de mejora
% ====================================================================

\chapter{Tema 5: Selección e Integración}

% ---- Bloque Julián ----
% ====================================================================
% TEMA 5 — JULIÁN
% Etapas y Herramientas Iniciales de Selección
% ====================================================================

\section{Tema 5: Etapas y Herramientas Iniciales de Selección}

\subsection{Etapas del proceso de selección}

\subsubsection{Preselección}

\subsubsection{Pruebas de selección}

\subsubsection{Decisión y contratación}

\subsection{Herramientas de selección básicas}

\subsubsection{Pruebas escritas}

\subsubsection{Pruebas orales}

\subsection{Gestión de entrevistas}

\subsubsection{Entrevistas individuales}

\subsubsection{Entrevistas colectivas}


% ---- Bloque Ismael ----
% ====================================================================
% TEMA 5 — ISMAEL
% Herramientas Avanzadas y Tendencias de Selección
% ====================================================================

\section{Tema 5: Herramientas Avanzadas y Tendencias de Selección}

\subsection{Técnicas de simulación en selección}

\subsection{Uso de assessment centers}

\subsection{Nuevas tendencias en selección}

\subsubsection{Vídeo entrevistas}

\subsubsection{Digitalización del proceso de selección}

\subsubsection{Retos y oportunidades de implantación}


% ---- Bloque David ----
% ====================================================================
% TEMA 5 — DAVID
% Acogida, Integración y Propuestas Finales
% ====================================================================

\section{Tema 5: Acogida, Integración y Propuestas Finales}

\subsection{Procesos de socialización en la organización}

\subsubsection{Socialización formal}

\subsubsection{Socialización informal}

\subsection{Manual de acogida: estructura y utilidad}

\subsection{Propuestas de mejora en selección e integración}

\subsubsection{Mejoras en la fase de incorporación}

\subsubsection{Mejoras de seguimiento durante la integración}

\subsubsection{Indicadores de éxito propuestos}



% -------------------------------------------------------------------
%  BACKMATTER
%  (bibliografía y fuentes consultadas)
% -------------------------------------------------------------------
\backmatter
\pdfbookmark[-1]{Referencias}{BM-Referencias}

% Fuentes consultadas (URLs y documentos teóricos)
% ====================================================================
%  BIBLIOGRAFÍA Y FUENTES CONSULTADAS
% ====================================================================

\chapter*{Fuentes consultadas}
\addcontentsline{toc}{chapter}{Fuentes consultadas}

A continuación se recogen las principales fuentes consultadas para la
elaboración del análisis práctico de los procesos de afectación (Temas 4 y 5) 
aplicados al caso de CaixaBank.

\begin{enumerate}[label={[\arabic*]}]

  \item \textbf{Guión básico para el desarrollo del trabajo de los Temas 4 y 5},
        Dirección de Recursos Humanos~I, Doble Grado ADE-Ingenierías,
        Universidad de Granada.

  \item \textbf{Marco Teórico de Recursos Humanos (Temas 4 y 5)} --- 
        Material docente y apuntes teóricos correspondientes a los procesos de 
        Reclutamiento, Selección y Socialización de la asignatura Dirección 
        de Recursos Humanos I.

  \item \textbf{Entrevistas prácticas (Fuentes primarias)} --- 
        Entrevistas cualitativas realizadas por los integrantes del grupo 
        a empleados actuales de CaixaBank. Estas entrevistas han servido como base 
        empírica exclusiva para analizar la aplicación real de las políticas de 
        reclutamiento, selección, integración y nuevas tendencias tecnológicas en la entidad.

\end{enumerate}

\endinput

% Bibliografía BibTeX (si se añaden citas con \cite{})
% \bibliographystyle{alpha-es}
% \begin{small}
%   \bibliography{library.bib}
% \end{small}

% !TeX root = ../tfg.tex
% !TeX encoding = utf8
%
% ====================================================================
%  APÉNDICE A — PREGUNTAS REALIZADAS A LA EMPRESA Y RESPUESTAS
%
%  RECORDATORIO AL EQUIPO:
%  - Este apéndice es OBLIGATORIO según las instrucciones de la asignatura.
%  - Debéis incluir las preguntas que realizasteis en la entrevista/visita
%    a la empresa y las respuestas que os dieron.
%  - Organizad las preguntas por temática (Tema 2 y Tema 3) o por el
%    orden en que se realizaron.
%  - Podéis incluir el nombre o cargo del entrevistado (si da permiso)
%    y la fecha y medio de la entrevista (presencial, videollamada...).
%  - Si hay información que el entrevistado pidió que no se publicara,
%    indicadlo (p.ej. "Dato confidencial — no se puede reproducir").
%  - Si grabasteis la entrevista (con permiso), podéis transcribir los
%    fragmentos más relevantes.
% ====================================================================

\chapter{Entrevista a CaixaBank: preguntas y respuestas}

% ============================================================
\section{Datos de la entrevista}
% ============================================================

\begin{table}[H]
  \centering
  \begin{tabularx}{\textwidth}{lX}
    \toprule
    \textbf{Parámetro}       & \textbf{Detalle} \\
    \midrule
    Empresa                  & CaixaBank \\
    Nombre del entrevistado  & [Anónimo a petición] \\
    Cargo                    & [Director de Recursos Humanos regional] \\
    Departamento             & [Recursos Humanos] \\
    Fecha                    & [12/04/2026] \\
    Duración                 & [15 min] \\
    Formato                  & [Telefónica] \\
    \bottomrule
  \end{tabularx}
\end{table}

% ============================================================
\section{Preguntas sobre el Tema 4 — Proceso de reclutamiento}
% ============================================================

\subsubsection*{¿Cuáles son los costes directos e indirectos asociados habitualmente al proceso de reclutamiento en la empresa? ¿Existe un presupuesto limitado para esto?}

\textbf{Respuesta:}\newline
``(El entrevistado sólo supo que acudían a una ETT para parte de este proceso).''

\subsubsection*{¿Cuál es el tiempo medio estimado (Time to Fill) desde que se detecta la necesidad de cubrir un puesto hasta que el candidato se incorpora?}

\textbf{Respuesta:}\newline
``(No supo responder el entrevistado).''

\subsubsection*{¿Qué enfoque de reclutamiento predomina actualmente en la organización: interno, externo o mixto?}

\textbf{Respuesta:}\newline
``Mixto. Normalmente el proceso de una vacante comienza en un portal interno en el cuál prevalece por un periodo de tiempo estipulado para cada tipo de vacante, de forma que predomina el ``moving'' para personal interno, ofreciendo así la posibilidad a los compañeros de adquirir nuevas responsabilidades y cambiar su ámbito de trabajo a personas que realmente cumplan los requisitos para la misma. También se valoran diferentes aspectos del proyecto en el que la persona esté colaborando actualmente para conocer si es posible la reducción de ``recursos'' asignados, se valorarán las cualidades del empleado para llevar a cabo un estudio de la aportación de conocimientos del ámbito, soft skills y hard skills, para elaborar una puntuación sobre la eficiencia de la persona en ambos proyectos para así poder optar por un bien estar para el empleado, además de sacar partido a su ``expertise''.''

\subsubsection*{A nivel interno: ¿Qué canales o métodos se utilizan para comunicar las vacantes a los empleados actuales (intranet, tablón de anuncios, programas de referidos)?}

\textbf{Respuesta:}\newline
``A nivel interno la forma de comunicar las vacantes disponibles a partir de la propia intranet ya mencionada.''

\subsubsection*{A nivel externo: ¿Cuáles son las fuentes de reclutamiento externo más efectivas para la empresa (portales de empleo, agencias de colocación, universidades, ferias de empleo)?}

\textbf{Respuesta:}\newline
``Dependiendo de la campaña utilizan diferentes medios de publicación, adaptando los canales de comunicación al tipo de perfil que buscan y al objetivo estratégico. Por ejemplo, si se trata de atraer talento joven para programas específicos como el CaixaBank Talent Program, es habitual reforzar la presencia en universidades, ferias de empleo y campañas digitales segmentadas en redes sociales profesionales.

En cambio, para perfiles más senior o especializados, se priorizan portales como LinkedIn o colaboraciones con consultoras de selección y headhunters. Asimismo, plataformas como InfoJobs o Glassdoor permiten dar visibilidad masiva a las vacantes.

De esta forma, la estrategia de reclutamiento externo no es única, sino que se ajusta en función del tipo de campaña, el público objetivo y la urgencia o especialización del perfil requerido, combinando distintos canales para maximizar el alcance y la calidad de las candidaturas.''

\subsubsection*{¿Se usan nuevas tendencias de reclutamiento, como por ejemplo 2.0 (Linkedin, Twitter, Facebook para atraer candidatos), 3.0 (employer branding) o video cv?}

\textbf{Respuesta:}\newline
``Actualmente, la mayoría de empresas utilizan una combinación de herramientas automatizadas e inteligencia artificial para filtrar grandes volúmenes de candidaturas de forma rápida y eficiente. 

La técnica más utilizada es el uso de sistemas ATS (Applicant Tracking System), que actúan como un primer filtro automatizado. Estos sistemas analizan los currículums, extraen la información y clasifican a los candidatos en función de su adecuación al puesto. También se usa machine learning para comparar a los candidatos con las ofertas a las que aplican y predecir su encaje, además de otras aplicaciones de la IA como análisis predictivos o automatización de distintos procesos como programar entrevistas.

Video cv en principio no, pero sí aceptan cartas de recomendación.''

\subsubsection*{¿Qué nuevas estrategias, tecnologías o tendencias tiene en mente aplicar la empresa en un futuro próximo para optimizar la captación de talento?}

\textbf{Respuesta:}\newline
``(No se le preguntó al entrevistado.)''

% ============================================================
\section{Preguntas sobre el Tema 5 — Proceso de selección e integración}
% ============================================================

\subsubsection*{Preselección: ¿Qué criterios (filtros killer, palabras clave, experiencia) se utilizan para realizar la criba curricular inicial?}

\textbf{Respuesta:}\newline
``Para la criba inicial se hace uso de sistemas ATS (Applicant Tracking System), que tienen unos filtros automatizados para cribar currículums. Estos filtros se basan en palabras clave, y se tiene en cuenta tanto filtros ``killer'' (requisitos excluyentes como nivel de idioma o formación específica dependiendo el puesto), experiencia profesional y coherencia del CV (progresión, estabilidad, adecuación al puesto), siendo este el criterio más importante, y la ubicación geográfica y disponibilidad.''

\subsubsection*{Pruebas: ¿Cómo se determina qué tipo de pruebas de selección se aplicarán a cada perfil profesional?}

\textbf{Respuesta:}\newline
``Dependiendo del tipo de perfil profesional o del nivel de responsabilidad del puesto. A perfiles técnicos se les suele plantear un caso práctico o real, mientras que a perfiles junior más pruebas de razonamiento para ver su potencial.''

\subsubsection*{Decisión y contratación: ¿Quiénes intervienen en la toma de decisión final (RRHH, responsable de departamento) y cómo se comunica la oferta al candidato elegido?}

\textbf{Respuesta:}\newline
``Intervienen el departamento de RRHH junto con los directores de área. Comunica la decisión recursos humanos para puestos bajos, o el director y oficial de RRHH para puestos más altos.''

\subsubsection*{¿En qué casos específicos se aplican pruebas escritas u orales (psicotécnicos, conocimientos, idiomas) y cómo se evalúan?}

\textbf{Respuesta:}\newline
``(El entrevistado sabe de la realización de pruebas escritas psicotécnicas en el pasado, pero desconoce qué ocurre en la realidad).''

\subsubsection*{¿Qué estructura siguen las entrevistas individuales (estructuradas, semi-estructuradas, por competencias) y qué rol juegan las entrevistas colectivas o dinámicas de grupo?}

\textbf{Respuesta:}\newline
``Son entrevistas estructuradas y se realizan dinámicas de grupo en ocasiones.''

\subsubsection*{¿Se utilizan técnicas de simulación (Role-play, método del caso) o Centros de Evaluación (Assessment Centers) para puestos directivos o clave?}

\textbf{Respuesta:}\newline
``Sí se usan.''

\subsubsection*{¿Se han implementado las vídeo entrevistas (ya sean en directo vía videollamada o en diferido/grabadas) en las primeras fases del proceso?}

\textbf{Respuesta:}\newline
``Se realizan video entrevistas, pero no suelen realizarse en las primeras fases del proceso, sino más adelante.''

\subsubsection*{Formal: ¿Qué acciones programadas (cursos de formación, reuniones de bienvenida, presentación de normativas) componen la socialización del nuevo empleado?}

\textbf{Respuesta:}\newline
``Sí hay, tanto para nuevos como antiguos. Tanto presenciales como online, cursos, reuniones, programas de bienvenida, una formación inicial obligatoria, etc.''

\subsubsection*{Informal: ¿Cómo se fomenta la integración social y cultural del empleado en el día a día (desayunos, asignación de un ``compañero'' o mentor no oficial, eventos de teambuilding)?}

\textbf{Respuesta:}\newline
``Se realizan tanto eventos sociales, como desayunos, se utilizan las mentorías para nuevos empleados, y eventos de teambuilding.''

\subsubsection*{¿Dispone la organización de un Manual de Acogida estructurado y actualizado?}

\textbf{Respuesta:}\newline
``Sí existe, mencionando principalmente los valores de la entidad, así como las funciones concretas del puesto. Sí está actualizado.''

\subsubsection*{¿Qué métricas o feedback se recogen de los nuevos empleados para evaluar la eficacia del proceso de selección y acogida?}

\textbf{Respuesta:}\newline
``Se realizan evaluaciones periódicas (6 meses, 1 año) mencionando qué cualidades destacan para el individuo en el puesto, como aquellas que son mejorables.''

\subsubsection*{¿Qué innovaciones tecnológicas (ej. Inteligencia Artificial para la preselección, gamificación en pruebas) o cambios metodológicos planea incorporar la empresa próximamente en esta área?}

\textbf{Respuesta:}\newline
``IA sí, gamificación también se está implementando.''



\end{document}
